\documentclass[11pt, a4paper]{article}

\usepackage[utf-8]{inputenc}
\usepackage[margin=1in]{geometry}
\usepackage{graphicx}
\usepackage{amsmath}
\usepackage{amssymb}
\usepackage{booktabs}
\usepackage{hyperref}
\usepackage{xcolor}
\usepackage{listings}
\usepackage{textcomp}

\title{\textbf{Bird Sound Classification using Transformers} \\
       Fine-Tuning Audio Spectrogram Transformers on BirdCLEF 2021}
\author{}
\date{}

\begin{document}

\maketitle

\section{Problem Formulation}

\subsection{Task Description}
This project addresses automated bird species classification from audio recordings. Given a raw audio clip (8--30 seconds), the model must predict which of 10 bird species produced the sound. This is a multi-class classification problem with significant real-world applications in ecological monitoring, biodiversity research, and citizen science initiatives.

\subsection{Dataset}
We use a curated subset of the BirdCLEF 2021 Kaggle competition dataset containing 10 acoustically diverse bird species:
\begin{itemize}
  \item \textbf{Total recordings:} 5,000 (500 per species from XenoCanto)
  \item \textbf{Audio format:} `.ogg` files at 16 kHz sample rate
  \item \textbf{Duration:} 8--30 seconds per recording
  \item \textbf{Data split:} Stratified into train (4,000 recordings / 400 per species), validation (500 / 50 per species), and test (500 / 50 per species)
\end{itemize}

\subsection{Data Preprocessing}
All audio is resampled to 16 kHz and chunked into non-overlapping 10-second segments (160,000 samples per chunk). Chunks shorter than 10 seconds are zero-padded. This produces approximately 24,000 training chunks, 3,000 validation chunks, and 3,000 test chunks. Recording-level predictions aggregate logits over all chunks of a full recording via mean pooling, reflecting real-world inference where we classify full-length clips.

\subsection{Evaluation Metrics}
We report two accuracy metrics:
\begin{itemize}
  \item \textbf{Chunk-level accuracy:} Fraction of 10-second chunks classified correctly (one prediction per chunk).
  \item \textbf{Recording-level accuracy:} Fraction of full recordings classified correctly after averaging chunk logits and taking argmax.
\end{itemize}
Recording-level accuracy is the primary metric because it reflects how the model would be deployed in practice.

\section{Model Design Choices}

\subsection{Architecture: Audio Spectrogram Transformer}
We selected the \textbf{Audio Spectrogram Transformer (AST)} as our core model. AST is a vision transformer pretrained on AudioSet (1.9M audio clips across 527 label categories). The architecture operates on 128-bin mel-spectrograms (1024-sample window, 128-sample hop) processed by a patch embedding layer followed by a 12-layer transformer encoder with 12 attention heads.

\textbf{Rationale:}
\begin{enumerate}
  \item \textbf{Transfer Learning:} AudioSet pretraining provides rich acoustic representations that generalize well to novel tasks and smaller datasets.
  \item \textbf{Attention Mechanism:} Transformers capture long-range temporal and spectral dependencies superior to CNNs for audio, important for discriminating bird vocalizations that encode species identity across time.
  \item \textbf{Efficiency:} At ~85M parameters, AST balances model capacity with computational tractability on moderate GPUs.
\end{enumerate}

\subsection{Fine-Tuning Strategy}
Rather than training from scratch, we fine-tune the pretrained AST backbone. This approach:
\begin{itemize}
  \item Dramatically reduces training time (5 epochs vs. hundreds on random initialization).
  \item Requires minimal labeled data (4,000 training recordings is small for deep learning).
  \item Achieves superior generalization by leveraging broad acoustic knowledge from AudioSet.
\end{itemize}

\noindent
The fine-tuning objective replaces the 527-class AudioSet head with a 10-class linear classifier and uses cross-entropy loss. All transformer parameters are unfrozen to allow adaptation to bird calls.

\section{Training Strategy and Hyperparameter Tuning}

\subsection{Hyperparameter Search}
Given the high computational cost of each full 5-epoch fine-tune (~65 minutes on an L40S GPU), we performed a \textbf{random search} to explore the learning rate $(lr)$ and weight decay $(wd)$ parameter space. The search spanned:
\begin{itemize}
  \item Learning rate: $[10^{-5}, 10^{-3}]$ (log-uniform)
  \item Weight decay: $[10^{-4}, 10^{-1}]$ (log-uniform)
\end{itemize}

\noindent
We evaluated $\approx 20$ configurations and selected the three best (spanning different regimes) for final validation and test evaluation. These three configurations are:

\begin{center}
\begin{tabular}{lccl}
  \toprule
  \textbf{Config} & \textbf{LR} & \textbf{WD} & \textbf{Description} \\
  \midrule
  A & $2.10 \times 10^{-4}$ & $5.70 \times 10^{-4}$ & Sweep best; standard regime \\
  B & $1.72 \times 10^{-4}$ & $7.41 \times 10^{-2}$ & Mid-LR, heavy regularization \\
  C & $1.30 \times 10^{-5}$ & $5.09 \times 10^{-2}$ & Low-LR, heavy regularization \\
  \bottomrule
\end{tabular}
\end{center}

\subsection{Training Hyperparameters}
All three configurations use:
\begin{itemize}
  \item \textbf{Optimizer:} AdamW
  \item \textbf{Batch size:} 16 (with gradient accumulation $\times 2$ $\Rightarrow$ effective batch size 32)
  \item \textbf{Epochs:} 5
  \item \textbf{Warmup steps:} 50
  \item \textbf{Loss:} Cross-entropy
  \item \textbf{Training data:} 100\% of training split (24k chunks)
  \item \textbf{Hardware:} Modal L40S GPU
\end{itemize}

\noindent
The three selected configurations demonstrate different regularization philosophies:
\begin{itemize}
  \item \textbf{Config A:} Light regularization, closest to the hyperparameter sweep optimum.
  \item \textbf{Config B \& C:} Heavy weight decay ($> 0.05$), reducing risk of overfitting but at risk of underfitting if too strong.
\end{itemize}

\noindent
Training dynamics show all three configurations converge to $\geq 0.91$ chunk-level validation accuracy by epoch 1, indicating the AudioSet pretrained backbone already captures discriminative acoustic features. Subsequent epochs provide modest refinement. Config C (low LR) exhibits the smoothest training curves, while configs B and C maintain visibly higher training loss than A due to regularization.

\section{Performance Analysis}

\subsection{Summary Results}
\begin{center}
\begin{tabular}{lccc}
  \toprule
  \textbf{Configuration} & \textbf{Best Val Acc} & \textbf{Test Chunk Acc} & \textbf{Test Record Acc} \\
  \midrule
  A & 0.9481 & 0.9468 & 0.9640 \\
  B & 0.9356 & 0.9279 & 0.9560 \\
  C & 0.9426 & 0.9362 & 0.9600 \\
  \bottomrule
\end{tabular}
\end{center}

\noindent
\textbf{Key observation:} Recording-level accuracy exceeds chunk-level by $\approx 1.5\%$. This improvement arises from mean-logit aggregation smoothing out occasional misclassifications on individual noisy chunks---a realistic benefit when classifying full-length clips at inference.

\subsection{Per-Class Performance}
F1 scores vary across species (range: $0.92$--$0.98$ at chunk level), with no systematic failures. Acoustic similarity between species (e.g., corvids, warblers) results in confusion, but all models consistently recover these patterns, indicating dataset-level challenges rather than configuration artifacts.

\subsection{Model Robustness}
The three configurations show high correlation in per-class F1 scores, suggesting the hyperparameter landscape near the search optimum is stable. Even config C (aggressively regularized) maintains strong performance, demonstrating that the pretrained AST foundation is robust across reasonable regularization choices.

\section{Limitations and Future Work}

\subsection{Limitations}
\begin{enumerate}
  \item \textbf{Dataset size:} 5,000 recordings (500 per species) is modest by deep learning standards. More data would likely further improve performance.
  \item \textbf{Recording length heterogeneity:} The dataset contains clips ranging from 8--30 seconds. Longer clips are chunked, potentially losing temporal context at chunk boundaries. Overlap-based or hierarchical chunking might help.
  \item \textbf{Acoustic diversity:} Each species is represented by recordings from XenoCanto, introducing author/equipment/recording-condition variation. Out-of-distribution audio (e.g., bioacoustic sensors, live field recordings with interference) may degrade performance.
  \item \textbf{Hyperparameter search scope:} Our random search covered $\approx 20$ configurations. A more extensive Bayesian optimization or grid search may identify further improvements, though compute constraints limit this.
  \item \textbf{Ensemble methods unexplored:} Averaging predictions from multiple models (AST, ResNet, other architectures) could boost accuracy, but was not explored due to time constraints.
\end{enumerate}

\subsection{Future Directions}
\begin{itemize}
  \item \textbf{Multi-modal learning:} Incorporate spectrogram features and raw waveform analysis jointly.
  \item \textbf{Attention visualization:} Analyze which time-frequency regions AST focuses on for each species to gain ecological insight.
  \item \textbf{Few-shot learning:} Evaluate model performance on rare species or species with limited training data.
  \item \textbf{Streaming inference:} Adapt the model for real-time processing suitable for autonomous recording units deployed in the field.
  \item \textbf{Domain adaptation:} Fine-tune on geographically or seasonally distinct subsets to assess generalization.
\end{itemize}

\section{Conclusion}

We successfully fine-tuned the Audio Spectrogram Transformer on 10 bird species from BirdCLEF 2021, achieving $\approx 96\%$ recording-level test accuracy. The pretrained AudioSet backbone combined with transfer learning enabled rapid convergence (within 1 epoch). Systematic hyperparameter exploration identified three stable configurations spanning different regularization regimes, all achieving comparable performance. The model generalizes well despite modest dataset size, underscoring the value of large-scale pretraining for audio tasks. Future work will explore ensemble methods, multi-modal representations, and domain adaptation for improved robustness in field deployment.

\end{document}
